\documentclass[11pt]{scrartcl}
\usepackage[sexy]{evan} %BLBLBLBLBLBLBLBLBLBLBLB EVANCHEN
\usepackage[T1]{fontenc}
\usepackage[utf8]{inputenc}
\usepackage{graphicx}
\usepackage{amsmath}
\usepackage{tikz}
\usepackage{tikz-cd}
\usepackage{asymptote}
\usepackage{amsthm}
\usepackage{amssymb}
\usepackage{gensymb}
\usepackage{amsfonts}
\usepackage[english,italian]{babel}
\usepackage{quoting}
\quotingsetup{font=big}
\usepackage{booktabs}
\usepackage{caption}
\usepackage{lmodern}
\usepackage{faktor}
\usepackage{bbm} %oer fare gli insiemi indicatori
\usepackage{leftindex} %per gli ideali e gli indici a sinistra
\usepackage{hyperref}
\usepackage{epigraph} 
\usepackage{pgfplots}
\usepackage[export]{adjustbox}

%Pacchetti per i codici
\usepackage{listings}
\usepackage{xcolor}
\usepackage{algpseudocode}

%Analisi e probabilità
%----------------------------------
\newcommand{\dif}{\text{d}}
%%% the pure symbol, see https://tex.stackexchange.com/a/718194/4427
\newcommand{\cind}{\perp\mspace{-9mu}\perp}
%%% define \indep
\NewDocumentCommand{\indep}{e{_}}{%
	\mathrel{%
		\IfNoValueTF{#1}{% no subscript
			\cind
		}{% subscript
			\mathchoice{\mathop{\cind}\limits_{#1}}{\cind_{#1}}{\cind_{#1}}{\cind_{#1}}%
		}%
	}%
}
%-----------------------------------

%Algebra lineare e geometria
%----------------------------------
\newcommand{\tr}[1]{\text{tr}\!\left(#1\right)}
\renewcommand{\rank}[1]{\text{rank}\!\left(#1\right)}
\newcommand{\Immagine}[1]{\text{Im}\left(#1\right)}
\renewcommand{\ker}[1]{\text{Ker}\left(#1\right)}
\renewcommand{\AA}{\mathbb{A}}
\newcommand{\LL}{\mathbb{L}}
\newcommand{\SO}{\text{SO}}
%----------------------------------

%Fisica
%----------------------------------
\newcommand{\UDiMisura}[1]{\,\text{#1}}
\newcommand{\ihat}{\hat{\imath}}
\newcommand{\jhat}{\hat{\jmath}}
\newcommand{\khat}{\hat{k}}
\newcommand{\posit}{\vec{r}}
\newcommand{\veloc}{\dot{\vec{r}}}
\newcommand{\accel}{\ddot{\vec{r}}}
%----------------------------------

%Informatica
%----------------------------------
\definecolor{codegreen}{rgb}{0,0.6,0}
\definecolor{codegray}{rgb}{0.5,0.5,0.5}
\definecolor{codepurple}{rgb}{0.58,0,0.82}
\definecolor{backcolour}{rgb}{0.95,0.95,0.92}

\lstdefinestyle{overleafwiki}{
	backgroundcolor=\color{backcolour},   
	commentstyle=\color{codegreen},
	keywordstyle=\color{magenta},
	numberstyle=\tiny\color{codegray},
	stringstyle=\color{codepurple},
	basicstyle=\ttfamily\footnotesize,
	breakatwhitespace=false,         
	breaklines=true,                 
	captionpos=b,                    
	keepspaces=true,                 
	numbers=left,                    
	numbersep=5pt,                  
	showspaces=false,                
	showstringspaces=false,
	showtabs=false,                  
	tabsize=2
}

\lstset{style=overleafwiki}
\newcommand{\bigo}[1]{\mathcal{O}\!\left(#1\right)}
\newcommand{\bigomega}[1]{\Omega\!\left(#1\right)}
\newcommand{\bigtheta}[1]{\Theta\!\left(#1\right)}
%----------------------------------

%Insiemistica e Algebra
%----------------------------------
\newcommand{\eset}{\varnothing}
\newcommand{\parti}[1]{\mathcal{P}\!\left(#1\right)}
\newcommand{\inj}{\hookrightarrow}
\newcommand{\surj}{\twoheadrightarrow}
\newcommand{\bij}{\hookrightarrow\mathrel{\mspace{-15mu}}\rightarrow}
\newcommand{\ordine}[2]{\text{ord}_{#1}\left(#2\right)}
\newcommand{\normal}{\mathrel{\unlhd}}
\newcommand{\MCD}{\text{MCD}}
\renewcommand{\mcm}{\text{mcm}}
%----------------------------------


\begin{document}
	\title{Appunti di Algebra 2}
	\subtitle{Andrea Lucchini - UNIPD}
	\date{Settembre 2026 - Gennaio 2027}
	
	
	\maketitle
	
	\tableofcontents
	
	\newpage
	
	\section{Gruppi}
	
	Iniziamo togliendo un po' di ruggine da Algebra 1.
	
	\subsection{Gruppi ciclici}
	
	Un gruppo ciclico è un gruppo che ammette un generatore, ed è probabilmente la famiglia più semplice di gruppi che possiamo considerare. Iniziamo classificandoli:
	 
	\begin{theorem}[Classificazione dei gruppi ciclici]
		Ogni gruppo ciclico $G$ o è isomorfo a $\faktor{\ZZ}{n\ZZ}$ o a $\ZZ$.
	\end{theorem}
	\begin{proof}
		Dato un gruppo ciclico $G=\langle g \rangle$ possiamo costruire una mappa 
		\begin{align*}
			r:\ZZ & \to G \\
			z & \mapsto g^z
		\end{align*}
		che risulta essere un omomorfismo. La mappa è chiaramente suriettiva, quindi ci manca solo da capire quale sia il nucleo. Abbiamo alcuni casi:
		\begin{itemize}
			\item Se $\ker{r}=\left\{0\right\}$ abbiamo semplicemente che $G\cong \ZZ$.
			\item Se il kernel è non banale allora sicuramente se $0\ne z\in \ker{r}$ anche $-z\in \ker{r}$ (si vede facilmente notando che se $g^z=1$ allora $g^{-z}=1$). Questo vuol dire che nel kernel ci sarà sicuramente almeno un numero positivo, ovvero esisterà il minimo intero positivo $n$ contenuto in $\ker{r}$. \\
			Ma allora per la divisione euclidea per ogni altro $k\in\ker{r}$ possiamo fare la divisione euclidea
			$$k=nq+r$$
			da cui
			$$g^k=g^{nq+r}=g^r$$
			ma $n>r\ge 0$ quindi $r=0$, cioè $n$ divide in modulo tutti gli elementi del kernel. Ma questo vuol dire proprio che
			$$\ker{r}\cong n\ZZ$$
			cioè $G\cong \faktor{\ZZ}{n\ZZ}$ per il primo teorema d'omomorfismo.
		\end{itemize}
		In tutti i casi abbiamo raggiunto la tesi, quindi abbiamo finito.
	\end{proof}
	
	Cosa possiamo dire allora degli ordini? Per il caso $G\cong \ZZ$ abbiamo chiaramente che l'ordine di tutti gli elementi è infinito, ma nel caso finito abbiamo:
	\begin{proposition}[Ordini di potenze]
		In un gruppo ciclico $G=\langle g \rangle$ di ordine $n$ abbiamo che
		$$\ordine{G}{g^k}=\frac{n}{\MCD(n,k)}.$$
	\end{proposition}
	\begin{proof}
		Consideriamo un $t$ tale che
		$$g^{tk}=1$$
		allora 
		$$n\mid tk$$
		il che vuol dire che, se chiamiamo $d:= \MCD(n,k)$ e 
		\begin{gather*}
			n=dn' \\
			k=dk'
		\end{gather*}
		abbiamo quindi che 
		$$n'\mid tk'$$
		ma dato che $n'$ e $k'$ sono coprimi dobbiamo avere $n'\mid t$, il che dimostra proprio che $n'\mid \ordine{G}{g^k}$, da cui la tesi.
	\end{proof}
	Il prossimo passo è studiare i sottogruppi di $G$ ciclico:
	\begin{theorem}[Sottogruppi di un ciclico]
		Un gruppo ciclico isomorfo a $\ZZ$ ha i sottogruppi in biezione con i numeri naturali, mentre un gruppo ciclico di ordine $n$ ha i sottogruppi in biezione con i divisori di $n$.
	\end{theorem}
	\begin{proof}
		Possiamo di nuovo scrivere una funzione, essendo $S(G)$ il reticolo dei sottogruppi di $G$:
		\begin{align*}
			\alpha:\NN & \to S(G)\\
			n & \mapsto \langle g^n \rangle
		\end{align*}
		usando la divisione euclidea possiamo di nuovo mostrare che ogni sottogruppo $H\le G$ è ciclico, e quindi generato da un elemento della forma $g^n$. Questo ci basta per concludere che $\alpha$ è suriettiva, ma cosa possiamo dire dell'iniettività? Beh abbiamo due casi:
		\begin{itemize}
			\item Se $\abs{g}=\infty$ abbiamo che se
			$$\langle g^{m_1}\rangle=\langle g^{m_2}\rangle$$
			allora abbiamo che $m_1=m_2$ e quindi per ogni $a$ esiste un $b$ tale che
			$$g^{m_1a}=g^{m_2b}$$ 
			da cui, visto che l'ordine di $g$ è infinito
			$$m_1a=m_2b$$
			cioè $m_1\mid m_2 $ e $m_2\mid m_1$ e quindi abbiamo una biezione perfetta.
			\item Se $\abs{g}=n$ esisterà un minimo $m\in \NN$ tale che $g^m\in \langle g^k\rangle$ allora abbiamo che, dividendo
			$$k=mq+r$$
			da cui $r=0$, cioè abbiamo proprio che il numero di classi di equivalenza sono il numero di divisori di $n$, cioè abbiamo una biezione tra il numero di divisori di $n$ e il numero di sottogruppi.
		\end{itemize}
		In tutti i casi abbiamo raggiunto la tesi, e quindi abbiamo finito.
	\end{proof}
	L'ultima cosa che rimane da capire è che libertà abbiamo nello scegliere i generatori dei gruppi, l'identità che abbiamo trovato sopra
	$$\ordine{G}{g^k}=\frac{n}{(n,k)}$$
	ci dice che i generatori sono esattamente tanti quanti i $k$ coprimi con $n$, cioè $\varphi(n)$. Questo ci permette di dimostrare un risultato interessante:
	\begin{proposition}[Somma di Eulero]
		Per ogni $n\in\NN$ vale la seguente identità
		$$n=\sum_{d\mid n}\varphi(d).$$
	\end{proposition}
	\begin{proof}
		Prediamo un gruppo ciclico $G$ con $n$ elementi. \'E chiaro che ogni elemento ha uno e un solo ordine, quindi l'idea consiste nel raggruppare gli elementi per sottogruppo che generano, ovvero:
		$$n=\sum_{H\le G}\abs{\left\{g\in G:\langle g \rangle=H \right\} }$$
		ma dato che abbiamo dimostrato che i sottogruppi di $H$ sono in biezione con i divisori di $n$, e che il numero di generatori di un gruppo ciclico è esattamente uguale a $\varphi(\abs{H})$ possiamo scrivere semplicemente
		$$\sum_{H\le G}\abs{\left\{g\in G:\langle g \rangle=H \right\} }=\sum_{H\le G}\varphi\left(\abs{H}\right)=\sum_{d\mid n}\varphi(d)$$
		che era la tesi.
	\end{proof}
	E infine un esercizietto carino è:
	\begin{proposition}[Caratterizzazione dei ciclici con ordine primo]
		Un gruppo $G$ tale per cui ogni suo $H\le G$ soddisfa o $H=G$ o $H=\left\{1\right\}$ è ciclico di ordine primo.
	\end{proposition}
	\begin{proof}
		Prendiamo un elemento $g\in G$ diverso da $1$, allora dato che
		$$\langle g \rangle \le G$$
		e $\langle g \rangle \ne \left\{1\right\}$ dobbiamo avere che $\langle g \rangle=G$, cioè $G$ è ciclico ed è generato da tutti i suoi elementi. Infine se per assurdo l'ordine di $G$ non fosse primo avrebbe sottogruppi non banali (in quanto avrebbe divisori non banali), dunque $\abs{G}=p$ e abbiamo finito.
	\end{proof}
	
	\subsubsection{Esponenti}
	
	Cominciamo ora a introdurre qualche nozione nuova:
	
	\begin{definition}[Esponente di un gruppo]
		Dato un gruppo $G$ finito, definiamo il suo \textit{esponente} $\exp G$ come il più piccolo intero $n$ tale che 
		$$g^n=1$$
		per ogni $g\in G$.
	\end{definition}
	\`E facile osservare che 
	\begin{proposition}[Caratterizzazione dell'esponente per divisibilità]
		In un gruppo $G$ finito abbiamo sempre che
		$$\exp G= \underset{g\in G}{\mcm}(\ordine{G}{g})$$
	\end{proposition}
	\begin{proof}
		Segue facilmente dal fatto che 
		$$g^k=1 \iff \ordine{G}{g}\mid k$$
		per ogni $k\in \NN$.
	\end{proof}
	
	Il primo teorema importante che dimostriamo in questo corso è
	\begin{theorem*}[Caratterizzazione dei gruppi ciclici]
		Per ogni gruppo abeliano $G$ abbiamo che $G$ è ciclico se e solo se il suo esponente coincide con l'ordine, ovvero
		$$\exp G = \abs{G} \iff G \text{ è ciclico}.$$
	\end{theorem*}
	Per dimostrarlo abbiamo bisogno di un paio di utili lemmi:
	\begin{lemma}[Lemma dell'ordine]
		Sia $G$ finito, e siano $g_1,g_2\in G$ tali che:
		\begin{itemize}
			\item Commutano, ovvero $g_1g_2=g_2g_1$.
			\item Hanno ordine coprimo, ovvero $(\abs{g_1},\abs{g_2})=1$.
		\end{itemize}
		allora abbiamo che 
		$$\abs{g_1g_2}=\abs{g_1}\abs{g_2}.$$
	\end{lemma}
	\begin{proof}
		Chiamiamo questi due ordini $n_1$ ed $n_2$. Supponiamo allora di avere un $t$ tale che
		$$\left(g_1g_2\right)^t=1$$
		allora dato che i due elementi commutano possiamo distribuire la potenza, dunque
		$$g_1^tg_2^t=1$$
		ovvero
		$$g_1^t=g_2^{-t}$$
		ma questo vuol dire proprio che $g_1\in \langle g_2\rangle$ e allo stesso modo $g_2\in \langle g_1 \rangle$, cioè 
		$$g_1,g_2\in \langle g_1 \rangle \cap \langle g_2 \rangle$$
		ma dato che, per il teorema di Lagrange, l'intersezione deve avere ordine un divisore comune ai due gruppi, deve essere necessariamente banale. Dunque
		$$g_1^t=g_2^t=1$$
		da questo troviamo che $n_1\mid t$ cioè $n_2\mid t$ e dunque 
		$$n_1n_2\mid t$$
		d'altra parte è chiaro che $n_1n_2$ sia un esponente per cui $g_1g_2$ si annulla, e quindi abbiamo proprio la tesi.
	\end{proof}
	Un altro lemma dunque che ci servirà è:
	\begin{lemma}[Caratterizzazione dell'esponente per massimalità]
		Sia $G$ finito e abeliano, e consideriamo un elemento $g$ tale che ha ordine massimo. Allora abbiamo
		$$\exp G = \abs{g}.$$
	\end{lemma}
	\begin{proof}
		Scomponiamo l'ordine
		$$\abs{G}=p_1^{n_1}\dots p_t^{n_t}$$
		e diamo un nome anche all'ordine di $g$, scomponendolo
		$$m:=\abs{g}=p_1^{\alpha_1}\dots p_t^{\alpha_t}.$$
		Per dimostrare la tesi ci basta mostrare che per ogni $x\in G$ si ha $x^m=1$, ossia per ogni $x\in G$ abbiamo
		$$\ordine{G}{x}\mid m$$
		scomponiamo l'ordine di un generico $x$ allora
		$$\abs{x}=p_1^{\beta_1}\dots p_t^{\beta_t}$$
		dobbiamo provare che $\beta_i\le \alpha_i$ per ogni $i$. Supponiamo per assurdo allora che $\beta_1> \alpha_1$, allora abbiamo che
		$$\abs{g^{p_1^{\alpha_1}}}=p_2^{\alpha_2}\dots p_t^{\alpha_t}$$
		mentre per $x$ facciamo
		$$\abs{x^{p_2^{\beta_2}\dots p_t^{\beta_t}}}=p_1^{\beta_1}$$
		ma allora siamo nella condizione del lemma dell'ordine, dunque il prodotto ha ordine esattamente
		$$\abs{g^{p_1^{\alpha_1}}x^{p_2^{\beta_2}\dots p_t^{\beta_t}}}=p_1^{\beta_1}p_2^{\alpha_2}\dots p_t^{\alpha_t}$$
		che è assurdo, perchè avevamo supposto che $g$ fosse l'elemento con ordine massimo in tutto il gruppo! E dunque abbiamo finito.
	\end{proof}
	Siamo quindi finalmente pronti a dimostrare il teorema sopra, che riprendiamo:
	\begin{theorem}[Caratterizzazione dei gruppi ciclici]
		Per ogni gruppo abeliano $G$ abbiamo che $G$ è ciclico se e solo se il suo esponente coincide con l'ordine, ovvero
		$$\exp G = \abs{G} \iff G \text{ è ciclico}.$$
	\end{theorem}
	\begin{proof}
		Un verso è facile, infatti se 
		$$\exp G = \abs{G}$$
		allora per il lemma precedente esisterà un $g$ tale che 
		$$\abs{g}=\exp G$$
		cioè $g$ ha l'ordine del gruppo, cioè $G$ è ciclico. D'altra parte per quanto dimostrato sopra i sottogruppi dei ciclici abbiamo che anche 
	\end{proof}
	Un'applicazione divertente di questo è
	\begin{theorem}[Gruppi finiti di campi]
		Dato un campo $\FF$ consideriamo il suo gruppo delle unità $\FF^*$, se $G\le \FF^*$ è un gruppo finito, allora $G$ è ciclico.
	\end{theorem}
	\begin{proof}
		Per Lagrange abbiamo immediatamente che 
		$$e:=\exp G \le \abs{G}=:n$$
		consideriamo allora il polinomio a coefficienti in $\FF$ tale che
		$$p(x)=x^e-1$$
		allora per il teorema di Ruffini sappiamo che il numero di radici di $p$ è minore o uguale di $e$, ma dato che per ogni $g\in G$ abbiamo
		$$p(g)=g^e-1$$
		dobbiamo necessariamente avere $n\le e$, da cui per la caratterizzazione dei gruppi ciclici dimostrata sopra abbiamo la tesi.
	\end{proof}
	Vediamo ancora un esercizio:
	\begin{proposition}[I razionali sono infinitamente generati]
		Prendiamo il gruppo additivo $(\QQ,+)$ e $G$ un gruppo 
		$$G\le \left(\QQ,+\right)$$
		tale che sia finitamente generato. Allora $G$ è ciclico.
	\end{proposition}
	\begin{proof}
		Se $G$ è finitamente generato allora
		$$G=\left\langle \frac{a_1}{b_1},\dots,\frac{a_t}{b_t} \right\rangle$$
		ma dato che
		$$\frac{a_1}{b_1}=\frac{1}{b_1b_2\dots b_t}a_1b_2b_3\dots b_t$$
		e analogamente per le altre frazioni, dato che siamo in notazione additiva questo vuol dire che
		$$\frac{a_i}{b_i}\in \left\langle \frac{1}{b_1b_2\dots b_t} \right\rangle$$
		per ogni $i$, ovvero
		$$G\le \left\langle \frac{1}{b_1b_2\dots b_t} \right\rangle$$
		da cui la tesi.
	\end{proof}
	E appunto come corollario di ciò otteniamo che i razionali sono infinitamente generati, in quanto è molto facile mostrare che non sono ciclici.
	
	\subsection{Classi laterali e struttura quoziente}
	
	Continuiamo il ripasso introducendo la seguente nozione
	\begin{definition}[Classe laterale]
		Sia $H\le G$ un gruppo. Se introduciamo una relazione di equivalenza tale che
		$$x\sim y \iff y = xh \qquad h\in H$$
		allora le classi di equivalenza sono dette \textit{classi laterali sinistre}, e analogamente possiamo definire le \textit{classi laterali destre}, questi corrispondono con gli insiemi $xH$ e $Hx$. Possiamo inoltre definire un \textit{trasversale destro o sinistro} come un insieme di rappresentanti di tutti i laterali destri o sinistri in $G$. Denotiamo con $T$ un trasversale.
	\end{definition}
	Possiamo riassumere quanto appena detto con
	$$G=\bigsqcup_{t\in T}tH$$
	e da questo segue immediatamente il teorema di Lagrange, come già visto, ovvero
	$$\abs{G:H}=\abs{T}=\frac{\abs{G}}{\abs{H}}.$$
	Un modo carino per dire che lavorare a destra o lavorare a sinistra è perfettamente equivalente è il seguente: giriamo la formula di decomposizione trovata sopra prendendo gli inversi di tutti gli elementi:
	$$G=G^{-1}=\bigsqcup_{t\in T}\left(tH\right)=\bigsqcup_{t\in T}H^{-1}t^{-1}=\bigsqcup_{t\in T}Ht^{-1}$$
	cioè possiamo passare agilmente da laterali destri a laterali sinistri invertendo.
	\begin{proposition}[Trasversali]
		Se $K\le H \le G$ e $T$ è un trasversale di $H$ in $G$ e $U$ è un trasversale di $K$ in $H$ allora il trasversale di $K$ in $G$ è il prodotto dei trasversali.
	\end{proposition}
	\begin{proof}
		Per definizione di trasversale ogni $g\in G$ può essere scritto come
		$$g=th$$
		con $t\in T$ e $h\in H$. Allo stesso modo possiamo scrivere
		$$h=uk$$
		con $u\in U$ e $k\in K$. Dunque questo dimostra che ogni $g$ si scrive come
		$$g=tuk$$
		cioè
		$$G=TUK$$
		da questo allora ci basta dimostrare che $TU$ è effettivamente formato da elementi in classi di equivalenza distinte modulo $K$. D'altra parte se
		$$t_1u_1K=t_2u_2K$$
		allora 
		$$u_2^{-1}t_2^{-1}t_1u_1\in K$$
		ma per chiusura abbiamo
		$$t_2^{-1}t_1\in u_2Ku_1^{-1}\subseteq H$$
		ma questo, dato che $T$ è un trasversale di $H$ vuol dire che $t_1=t_2$, e quindi anche $u_1=u_2$ (dato che $U$ è un trasversale di $K$), da cui la tesi.
	\end{proof}
	Un'altra cosa che si può fare con gli indici e i trasversali ci viene suggerita dal seguente grafico 
	\begin{figure}[h]
		\begin{center}
			\begin{tikzcd}
				&&G\arrow[d]&& \\
				&&\langle H,K \rangle\arrow[dr]\arrow[dl] &&  \\
				&H\arrow[dr]&&K\arrow[dl]& \\
				&&H\cap K&&
			\end{tikzcd}
		\end{center}
	\end{figure}
	Notiamo che intuitivamente la distanza tra $G$ e $K$ è maggiore di quella tra $H$ e $H\cap K$, allora la disuguaglianza che vogliamo è
	\begin{proposition}[Disuguaglianza tra trasversali]
		Se $H\le G$ e $K \le G$ abbiamo che 
		$$\abs{H:H\cap K}\le \abs{G:K}$$
	\end{proposition}
	\begin{proof}
		Prendendo due laterali diversi di $H\cap K$ abbiamo
		$$h_1\left(H\cap K\right)\ne h_2\left(H\cap K\right)$$
		cioè $h_2^{-1}h_1 \not\in H\cap K$. Cioè $h_2^{-1}h_1\not\in K$ insomma abbiamo altre due classi laterali disgiunte per $K$, il che dimostra proprio la tesi.
	\end{proof}
	Da cui 
	\begin{corollary}[Disuguaglianza per l'intersezione degli indici]
		In generale se $G\ge H,K$ abbiamo
		$$\abs{G:H\cap K}\le \abs{G:H}\abs{G:K}.$$
	\end{corollary}	
	\begin{proof}
		Per quanto appena mostrato abbiamo
		\begin{gather*}
			\abs{G:K}\ge \abs{H:H\cap K} \\
			\abs{G:H}\ge \abs{K:H\cap K}
		\end{gather*}
		ovvero moltiplicando membro a membro e usando quanto mostrato sul prodotto dei trasversali abbiamo
		$$\abs{G:K}\abs{G:H}\ge  \abs{H:H\cap K} \abs{K:H\cap K}=\abs{HK}$$
	\end{proof}
	E purtroppo in generale non possiamo fare di meglio. La buona notizia è che:
	\begin{proposition}[Se gli indici sono coprimi si ha uguaglianza]
		SE GLI INDICI SONO COPRIMI ALLORA HAI UGUAGLIANZA, RIPRENDILA
	\end{proposition}
	Quando si parla di quozienti parliamo di gruppi normali:
	\begin{definition}[Gruppo Normale]
		Un sottogruppo $N\le G$ si dice normale e si scrive $N\normal G$ se e solo se 
		$$gHg^{-1}=H$$
		per ogni $g\in G$.
	\end{definition}
	Inoltre ricordiamo la definizione di prodotto diretto:
	\begin{definition}[Prodotto tra gruppi]
		Dati due $G\ge H,K$ allora definiamo il gruppo \textit{prodotto} come
		$$HK:=\left\{hk:h\in H, k\in K\right\}.$$
	\end{definition}
	Non è necessario che questo sia un gruppo! Infatti prendendo gli inversi di tutto troviamo
	$$\left(HK\right)^{-1}=K^{-1}H^{-1}=KH$$
	che vuol dire che se $HK$ è un sottogruppo (dunque chiuso rispetto agli inversi), dobbiamo necessariamente avere
	$$HK=KH$$
	ma questa condizione è anche sufficiente, infatti da questa possiamo anche dimostrare che $HK$ è chiuso rispetto al prodotto, ovvero
	$$h_1k_1h_2k_2=h_1h^*k^*k_2=(h_1h^*)(k^*k_2).$$
	Da questa discussione estraiamo quindi:
	\begin{corollary}[Normale implica prodotto un gruppo]
		Se $H\le G$ e $N\normal G$ abbiamo che $HK\le G$.
	\end{corollary}
	Con i prodotto possiamo fare tante cose interessanti:
	\begin{proposition}[Ordine del prodotto]
		Per ogni $H,K\le G$ abbiamo
		$$\abs{HK}=\frac{\abs{K}\abs{K}}{\abs{K\cap H}}.$$
	\end{proposition}
	\begin{proof}
		Consideriamo la funzione 
		$$\sigma : H\times K \to HK$$
		che per ogni coppia $(h,k)$ moltiplica gli elementi. Prima di tutto questa è sicuramente suriettiva, e il dominio ha proprio cardinalità $\abs{H}\abs{K}$. D'altra parte se abbiamo
		$$h_1k_1=h_2k_2$$
		allora
		$$h_2^{-1}h_1=k_2k_1^{-1}\in H\cap K$$
		cioè abbiamo che 
		$$(h_1,h_2)^{-1}(k_1,k_2)\in (H\cap K)^2 $$
		quindi se due elementi hanno la stessa immagine differiscono per un prodotto di $H\cap K$, d'altra parte è facile vedere come due elementi che differiscono per un prodotto per un elemento di $H\cap K$ hanno la stessa immagine, infatti se $x\in H\cap K$ abbiamo
		$$(hx,x^{-1}k)\mapsto hk$$
		dunque abbiamo proprio la tesi.
	\end{proof}
	Cosa succede allora se $N\normal G$, $H\le G$ e il prodotto è proprio $G=HN$, mentre $H\cap N=\left\{1\right\}$? Beh per la funzione costruita prima questo vuol dire che ogni elemento $g\in G$ possiamo trovare una e una sola coppia $(h,k)$ tale che
	$$hk=g$$
	quello che ci viene da pensare immediatamente allora è che la biezione $\sigma$ che abbiamo costruito sia un isomorfismo. E questo è \textbf{falso} in generale, infatti abbiamo che
	$$(h,1)\mapsto h$$
	e 
	$$(1,k)\mapsto k$$
	ma gli elementi a sinistra commutano (per come è fatto il prodotto diretto), mentre quelli a destra non necessariamente.
	\subsubsection{Applicazioni dei teoremi d'omomorfismo}
	Ricordiamo i teoremi d'omomorfismo: questi scaturiscono essenzialmente dalla proiezione canonica
	$$G\overset{\pi}{\longrightarrow}\faktor{G}{N}$$
	da questa equazione si ricava facilmente il primo teorema d'omomorfismo. Il secondo nasce osservando che questa funzione induce anche una mappa tra i reticoli, ovvero
	$$\pi_{\mid H}:H\to \faktor{HN}{N}$$
	e dato che 
	$$\ker{\pi_{\mid H}}=\ker{\pi}\cap H=N\cap H$$
	abbiamo proprio il secondo teorema d'omomorfismo. Più articolato è il terzo teorema: quando ci troviamo nella situazione $H,K\normal G$ possiamo costruire la mappa
	$$\gamma:\faktor{G}{K}\to \faktor{G}{H}$$
	tale che $gK\mapsto gH$. Da cui ricaviamo facilmente il terzo teorema d'omomorfismo. Un'applicazione interessante è:
	\begin{theorem*}[Cauchy dispari]
		Sia $G$ un gruppo tale che $\abs{G}=2m$, con $m$ dispari. Allora $G$ contiene un sottogruppo di ordine $m$.
	\end{theorem*}
	Per arrivarci tuttavia passiamo per un paio di lemmi:
	\begin{lemma}[Ordine $2$]
		Se $G$ è un gruppo finito di ordine pari allora contiene un elemento di ordine $2$.
	\end{lemma}
	\begin{proof}
		La dimostrazione consiste essenzialmente nel considerare la funzione che associa ad ogni $g\mapsto g^{-1}$ il suo inverso. Dato che il gruppo ha ordine pari, e che $(g^{-1})^{-1}$ abbiamo necessariamente che il numero degli elementi che vanno in se stessi sotto questa funzione deve necessariamente essere pari. D'altra parte questo insieme è sicuramente non vuoto, in quanto
		$$1\mapsto 1^{-1}=1$$
		e dunque dato che è pari ci sarà sempre almeno un altro elemento, che avrà ordine $2$ appunto.
	\end{proof}
	Un'altra cosa che ci servirà per dimostrare il teorema sopra è il teorema di Cayley. Ricordiamo che il teorema di Cayley scaturiva essenzialmente dal considerare la funzione
	\begin{align*}
		l_g:G &\to G \\
		x & \mapsto gx
	\end{align*}
	che è un omomorfismo. Non solo, è anche biettivo, dunque possiamo costruire una biezione $g\mapsto l_g$ che ci permette di identificare $G$ con $\text{Sym}(G)$, ovvero ogni gruppo finito è un gruppo di permutazioni. Finalmente quindi:
	\begin{theorem}[Cauchy dispari]
		Sia $G$ un gruppo tale che $\abs{G}=2m$, con $m$ dispari. Allora $G$ contiene un sottogruppo di ordine $m$.
	\end{theorem}
	\begin{proof}
		Consideriamo la funzione 
		\begin{align*}
			l:G& \to \text{Sym}(G) \\
			g &\mapsto l_g 
		\end{align*}
		la cosa cruciale è che ogni permutazione può essere scomposta in trasposizioni, ovvero esplicitamente in questo caso
		$$l(g)=(x_1\quad gx_1)(x_2 \quad gx_2)\dots (x_n \quad gx_n)$$
		e dato che l'ordine di $G$ è pari questo vuol dire che il numero di trasposizioni è $n-1$, dunque $l(g)$ non è nel gruppo alterno. Dato che $\text{Alt}(G)$ ha ordine $2$ deve essere normale, dunque
		$$\text{Alt}(G)\normal \text{Alt}(G)l(G)\normal \text{Sym}(G)$$
	\end{proof}
	
	\newpage
	
	\section{Soluzioni agli esercizi proposti}
	Durante il corso venivano proposti diversi esercizi nei fogli di riepilogo settimanali. Qui una raccolta di testi e soluzioni divisi per settimana.
	\subsection{Settimana 1: testi}
	
	\subsection{Settimana 1: soluzioni}
	
	\subsection{Settimana 2: testi}
	
	\subsection{Settimana 2: soluzioni}
	
	\subsection{Settimana 3: testi}
	
	\subsection{Settimana 3: soluzioni}
	
	\subsection{Settimana 4: testi}
	
	\subsection{Settimana 4: soluzioni}
	
	\subsection{Settimana 5: testi}
	
	\subsection{Settimana 5: soluzioni}
	
	\subsection{Settimana 6: testi}
	
	\subsection{Settimana 6: soluzioni}
	
\end{document}
